% ========================================================
% BÀI TẬP TOÁN 5 - KẾT NỐI TRI THỨC VỚI CUỘC SỐNG
% Bài 8: Ôn tập hình học và đo lường
% Định dạng: BTTeX Part III - Trắc nghiệm điền số ngắn
% ========================================================

% Câu 1
\begin{ex}
  Một góc bẹt bằng bao nhiêu góc vuông?
  \shortans{2}
  \loigiai{
    Một góc bẹt bằng $180^\circ$, một góc vuông bằng $90^\circ$. Vì $180 : 90 = 2$ nên góc bẹt bằng 2 góc vuông.
  }
\end{ex}

% Câu 2
\begin{ex}
  Trong một hình chữ nhật có bao nhiêu góc vuông?
  \shortans{4}
  \loigiai{
    Hình chữ nhật có 4 góc ở 4 đỉnh đều là góc vuông.
  }
\end{ex}

% Câu 3
\begin{ex}
  Đổi $3\text{ m } 4\text{ dm}$ thành bao nhiêu đề-xi-mét?
  \shortans{34}
  \loigiai{
    $3\text{ m} = 30\text{ dm}$. Vậy $3\text{ m } 4\text{ dm} = 34\text{ dm}$.
  }
\end{ex}

% Câu 4
\begin{ex}
  Đổi $5\text{ km } 20\text{ m}$ thành bao nhiêu mét?
  \shortans{5020}
  \loigiai{
    $5\text{ km} = 5\,000\text{ m}$. Vậy $5\text{ km } 20\text{ m} = 5\,020\text{ m}$.
  }
\end{ex}

% Câu 5
\begin{ex}
  Đổi $4\text{ tấn } 50\text{ kg}$ thành bao nhiêu ki-lô-gam?
  \shortans{4050}
  \loigiai{
    $4\text{ tấn} = 4\,000\text{ kg}$. Do đó $4\text{ tấn } 50\text{ kg} = 4\,050\text{ kg}$.
  }
\end{ex}

% Câu 6
\begin{ex}
  Đổi $7\text{ tạ } 5\text{ yến}$ thành bao nhiêu yến?
  \shortans{75}
  \loigiai{
    $7\text{ tạ} = 70\text{ yến}$. Do đó $7\text{ tạ } 5\text{ yến} = 75\text{ yến}$.
  }
\end{ex}

% Câu 7
\begin{ex}
  Đổi $6\text{ m}^2\text{ } 15\text{ dm}^2$ thành bao nhiêu đề-xi-mét vuông?
  \shortans{615}
  \loigiai{
    $1\text{ m}^2 = 100\text{ dm}^2$. Do đó $6\text{ m}^2\text{ } 15\text{ dm}^2 = 615\text{ dm}^2$.
  }
\end{ex}

% Câu 8
\begin{ex}
  Đổi $8\text{ dm}^2\text{ } 4\text{ cm}^2$ thành bao nhiêu xăng-ti-mét vuông?
  \shortans{804}
  \loigiai{
    $8\text{ dm}^2 = 800\text{ cm}^2$. Do đó $8\text{ dm}^2\text{ } 4\text{ cm}^2 = 804\text{ cm}^2$.
  }
\end{ex}

% Câu 9
\begin{ex}
  Một hình bình hành có độ dài đáy là $15\text{ cm}$ và chiều cao tương ứng là $8\text{ cm}$. Diện tích hình bình hành đó là bao nhiêu xăng-ti-mét vuông?
  \shortans{120}
  \loigiai{
    Diện tích hình bình hành là: $15 \times 8 = 120\text{ cm}^2$.
  }
\end{ex}

% Câu 10
\begin{ex}
  Một hình thoi có độ dài hai đường chéo lần lượt là $16\text{ cm}$ và $10\text{ cm}$. Diện tích hình thoi đó là bao nhiêu xăng-ti-mét vuông?
  \shortans{80}
  \loigiai{
    Diện tích hình thoi là: $(16 \times 10) : 2 = 80\text{ cm}^2$.
  }
\end{ex}

% Câu 11
\begin{ex}
  Một hình chữ nhật có chu vi $60\text{ cm}$, chiều dài gấp đôi chiều rộng. Diện tích hình chữ nhật đó là bao nhiêu xăng-ti-mét vuông?
  \shortans{200}
  \loigiai{
    Nửa chu vi là $60 : 2 = 30\text{ cm}$. Chiều rộng là $30 : (1 + 2) = 10\text{ cm}$. Chiều dài là $20\text{ cm}$. Diện tích là $20 \times 10 = 200\text{ cm}^2$.
  }
\end{ex}

% Câu 12
\begin{ex}
  Một mảnh đất hình vuông có chu vi là $64\text{ m}$. Diện tích mảnh đất đó là bao nhiêu mét vuông?
  \shortans{256}
  \loigiai{
    Cạnh hình vuông là $64 : 4 = 16\text{ m}$. Diện tích là $16 \times 16 = 256\text{ m}^2$.
  }
\end{ex}

% Câu 13
\begin{ex}
  Một thửa ruộng hình chữ nhật có chiều dài $40\text{ m}$, chiều rộng $25\text{ m}$. Người ta thu hoạch được $2\text{ kg}$ thóc trên mỗi mét vuông. Hỏi thửa ruộng thu hoạch được tất cả bao nhiêu tạ thóc?
  \shortans{20}
  \loigiai{
    Diện tích là $40 \times 25 = 1\,000\text{ m}^2$. Số thóc là $1\,000 \times 2 = 2\,000\text{ kg} = 20\text{ tạ}$.
  }
\end{ex}

% Câu 14
\begin{ex}
  Một căn phòng hình chữ nhật có chiều dài $8\text{ m}$, chiều rộng $6\text{ m}$. Người ta lát sàn bằng các viên gạch hình vuông cạnh $40\text{ cm}$. Cần dùng tất cả bao nhiêu viên gạch (bỏ qua mạch vữa)?
  \shortans{300}
  \loigiai{
    Diện tích phòng là $8 \times 6 = 48\text{ m}^2 = 480\,000\text{ cm}^2$. Diện tích 1 viên gạch là $40 \times 40 = 1\,600\text{ cm}^2$. Số gạch cần là: $480\,000 : 1\,600 = 300$ (viên).
  }
\end{ex}

% Câu 15
\begin{ex}
  Có bao nhiêu giây trong $3\text{ phút } 20\text{ giây}$?
  \shortans{200}
  \loigiai{
    $3\text{ phút} = 180\text{ giây}$. Tổng số giây là $180 + 20 = 200\text{ giây}$.
  }
\end{ex}

% Câu 16
\begin{ex}
  Năm $1010$ thuộc thế kỉ thứ mấy?
  \shortans{11}
  \loigiai{
    Năm 1010 thuộc thế kỉ thứ XI (thế kỉ 11).
  }
\end{ex}

% Câu 17
\begin{ex}
  Bác Hồ đọc Tuyên ngôn Độc lập năm $1945$. Năm đó thuộc thế kỉ thứ mấy?
  \shortans{20}
  \loigiai{
    Năm 1945 thuộc thế kỉ thứ XX (thế kỉ 20).
  }
\end{ex}

% Câu 18
\begin{ex}
  Một hình thoi có cạnh dài $15\text{ cm}$. Chu vi của hình thoi đó là bao nhiêu xăng-ti-mét?
  \shortans{60}
  \loigiai{
    Chu vi hình thoi là: $15 \times 4 = 60\text{ cm}$.
  }
\end{ex}

% Câu 19
\begin{ex}
  Người ta rào một khu vườn hình chữ nhật dài $18\text{ m}$, rộng $12\text{ m}$, có để một cửa ra vào rộng $2\text{ m}$. Chiều dài hàng rào là bao nhiêu mét?
  \shortans{58}
  \loigiai{
    Chu vi khu vườn là $(18 + 12) \times 2 = 60\text{ m}$. Chiều dài hàng rào là $60 - 2 = 58\text{ m}$.
  }
\end{ex}

% Câu 20
\begin{ex}
  Đổi $300\text{ dm}^2$ thành bao nhiêu mét vuông?
  \shortans{3}
  \loigiai{
    $1\text{ m}^2 = 100\text{ dm}^2$. Vậy $300\text{ dm}^2 = 3\text{ m}^2$.
  }
\end{ex}

% Câu 21
\begin{ex}
  Tam giác $ABC$ có một góc vuông và hai cạnh góc vuông lần lượt là $6\text{ cm}$ và $8\text{ cm}$. Diện tích tam giác đó là bao nhiêu xăng-ti-mét vuông?
  \shortans{24}
  \loigiai{
    Diện tích tam giác vuông là: $(6 \times 8) : 2 = 24\text{ cm}^2$.
  }
\end{ex}

% Câu 22
\begin{ex}
  Một hình bình hành có diện tích $72\text{ cm}^2$, độ dài đáy là $9\text{ cm}$. Chiều cao tương ứng của hình bình hành là bao nhiêu xăng-ti-mét?
  \shortans{8}
  \loigiai{
    Chiều cao là: $72 : 9 = 8\text{ cm}$.
  }
\end{ex}

% Câu 23
\begin{ex}
  Một cuộn dây thép dài $24\text{ m}$. Người ta cắt thành các đoạn dây nhỏ dài $8\text{ dm}$. Cắt được tất cả bao nhiêu đoạn dây?
  \shortans{30}
  \loigiai{
    $24\text{ m} = 240\text{ dm}$. Số đoạn dây là: $240 : 8 = 30$ (đoạn).
  }
\end{ex}

% Câu 24
\begin{ex}
  Một con chim bay với vận tốc $12\text{ m/s}$. Trong $15$ giây con chim đó bay được bao nhiêu mét?
  \shortans{180}
  \loigiai{
    Quãng đường bay được là: $12 \times 15 = 180\text{ m}$.
  }
\end{ex}

% Câu 25
\begin{ex}
  Nếu tăng cạnh của một hình vuông lên gấp đôi thì diện tích của nó tăng lên gấp mấy lần?
  \shortans{4}
  \loigiai{
    Diện tích hình vuông tỉ lệ với bình phương cạnh: $2 \times 2 = 4$. Vậy diện tích tăng 4 lần.
  }
\end{ex}
